% -----------------------------------------------------------------------------
% Section 6: Discussion
% Paper: "Silent Decisions Are Type Errors: Enforcing AI Accountability
%         via Linear Resource Types"
% Venue: IEEE Computer — Special Issue on AI Governance
% Deadline: 13 April 2026
% Authors: Daniel Lau Pereira Soares
% -----------------------------------------------------------------------------

\section{Discussion}
\label{sec:discussion}

The Constitutional Enclosure Theorem establishes a strong structural
guarantee within a precisely bounded scope.
This section examines the boundaries of that scope honestly,
identifies three limitations of the current reference implementation,
and outlines directions for future work.

% ----------------------------------------------------------------
\subsection{Scope and Boundaries of the Guarantee}
\label{sec:scope}

The theorem holds within Safe Rust and for code whose crate graph
contains \texttt{silent-decisions-proof}.
Two categories of code lie outside this perimeter.

\paragraph{Unsafe Rust.}
The \texttt{unsafe} keyword grants the programmer the ability to
bypass borrow-checker and visibility rules, including the private
field restrictions that enforce Protection~A
(Table~\ref{tab:protections}).
Production deployments that require the full strength of the theorem
should enforce a zero-\texttt{unsafe} policy over the transitive
dependency graph, auditable via \texttt{cargo geiger}.
This is not a theoretical limitation but an operational one:
the tooling exists, and its enforcement is a policy decision,
not a technical barrier.

\paragraph{Semantic truthfulness.}
The theorem guarantees that \emph{some} context was hashed and bound
to a verdict---not that the hashed context is the correct or complete
input presented to the AI system at inference time.
A malicious or inattentive caller could pass a fixed byte string
(\texttt{b"placeholder"}) to \texttt{EvidenceToken::new} and
obtain a structurally valid token.
Addressing semantic truthfulness requires attestation mechanisms at
the application layer (e.g., hardware-rooted measurement of the
inference pipeline), which is orthogonal to the type-level guarantee
we provide.

% ----------------------------------------------------------------
\subsection{Limitations of the Reference Implementation}
\label{sec:limitations}

The \texttt{silent-decisions-proof} kernel is a proof-of-concept
designed to validate the theoretical claims of this paper.
Three implementation choices trade production robustness for
clarity of exposition.

\paragraph{L1 — Static HMAC key.}
The current implementation uses a compile-time constant as the HMAC
signing key.
This is sufficient to demonstrate integrity verification (the HMAC
binds evidence, decision, and explanation together) but is
inadequate for production use, where key compromise would allow
post-hoc forgery of the integrity seal.
Production deployments should integrate a Hardware Security Module
(HSM) or a key-management service (e.g., AWS KMS, HashiCorp Vault)
for key issuance and rotation.

\paragraph{L2 — Public \texttt{ComplianceToken} constructor.}
The \texttt{ComplianceToken::new} constructor is \texttt{pub},
allowing any caller to self-declare a jurisdiction, policy version,
and contestability deadline without external validation.
In a production governance kernel, token issuance should be
restricted to a trusted policy engine that validates the declared
parameters against a verified regulatory database.
The reference implementation makes the constructor public to
accommodate the standalone binary demo; a production deployment would
restrict it to \texttt{pub(crate)} or require a signed token from
an authorised issuer.

\paragraph{L3 — Single-node evidence chain.}
The current \texttt{EvidenceToken} captures a single BLAKE3 hash
of the decision context.
Multi-step AI pipelines (retrieval, re-ranking, generation, scoring)
produce intermediate states that each carry accountability
implications.
Extending the type system to a \emph{chain} of evidence tokens---
where each pipeline stage consumes the prior token and produces
a new one---is a natural generalisation that preserves the
$V \multimap (E \otimes C)$ invariant at each step.

% ----------------------------------------------------------------
\subsection{The CAL Design Principle}
\label{sec:cal}

The three limitations above reflect a deliberate design choice that
we call the \emph{CAL principle}: in a governance system under
resource constraints, the three properties of
\textbf{C}onsistency (every verdict has evidence),
\textbf{A}vailability (the system can always issue a verdict), and
\textbf{L}egality (tokens are issued only by authorised entities)
cannot all be maximised simultaneously.
The reference implementation prioritises Consistency and Availability
over Legality---a reasonable choice for a proof-of-concept, but
one that must be revisited before production deployment.
The CAL principle provides a vocabulary for communicating these
trade-offs to system architects and regulators.

% ----------------------------------------------------------------
\subsection{Future Work}
\label{sec:future}

Three directions are most immediate.
First, integrating HSM-backed key management (addressing L1) and a
regulatory-oracle-signed \texttt{ComplianceToken} (addressing L2)
would close the gap between the reference implementation and a
production governance kernel.
Second, extending the evidence chain to multi-step pipelines
(addressing L3) would make the framework applicable to
retrieval-augmented generation and agentic AI architectures.
Third, formalising the type-system extension in a proof assistant
(e.g., Lean~4 or Coq) would elevate the Constitutional Enclosure
Theorem from a machine-checked Rust proof to a fully mechanised
mathematical proof, strengthening its standing as a foundational
result in the theory of accountable computation.
